% Pista ana-23j-q3 · Anàlisi
% Procedència: PAU Catalunya 2023 · Convocatòria de juny · Sèrie 1 · Qüestió 3
\documentclass[11pt,a4paper]{article}
\usepackage{pista}
\pistainfo{Catalunya 2023}{Convocatòria de juny}{Sèrie 1}

\begin{document}

\section*{\textcolor{blauPista}{Qüestió 3}}

\noindent Sigui
\[
f'(x) = \begin{cases} x - 1, & \text{si } x \le 2 \\[2pt] \dfrac{1}{x - 1}, & \text{si } x > 2 \end{cases}
\]
la funció derivada d'una funció derivable $f(x)$ que passa pel punt $A = (0, 3)$.

\subsection*{\textit{a)} \normalfont Calculeu la funció $f(x)$. \hfill\textnormal{[1,5 punts]}}

\pista{Per obtenir l'expressió de $f(x)$ a partir de les dades, hauràs d'integrar cada tram per separat. Cada primitiva porta una constant d'integració arbitrària, obtindràs \emph{una per a cada tram}, és a dir, $C_{1}$ i $C_{2}$. Per fixar-les, tens dues condicions: $f(0) = 3$ et determina la constant del primer tram, i com que $f$ és derivable a tot arreu (i, per tant, també contínua), la continuïtat a $x = 2$ et determina la constant del segon tram.}

\subsection*{\textit{b)} \normalfont Calculeu l'equació de la recta tangent a la funció $f'(x)$ en el punt d'abscissa $x = 3$. \hfill\textnormal{[1 punt]}}

\pista{Atenció a l'enunciat: et demanen la recta tangent a $f'(x)$, \emph{no} a $f(x)$. La fórmula de la recta tangent a una funció $g$ en el punt d'abscissa $a$ és $y - g(a) = g'(a)(x - a)$: aquí, la funció $g$ és $f'$, així que la seva derivada és $f''$. Com que $x = 3$ cau dins del segon tram (és a dir, $x > 2$), utilitza la fórmula de $f'(x)$ corresponent a aquest tram, deriva-la per obtenir $f''(3)$ i avalua $f'(3)$. Amb aquests dos valors, obtenir l'equació de la recta que et demanen és senzill.}

\end{document}
